\section{ماژول zerokv\_neo\_compressor (ZeroKV-Neo Compressor)}
\label{sec:module_zerokv_neo_compressor}

\subsection{معرفی و انگیزه}
این ماژول در چارچوب پژوهش حاضر پیاده‌سازی و ارزیابی شده و فاقد مقاله مستقل جداگانه در ادبیات تحقیق است.

در مدل‌های زبانی خودبازگشت، رشد خطی اندازه حافظه نهان کلید-مقدار (\lr{KV Cache}) به موازات افزایش طول توالی ورودی، مهم‌ترین گلوگاه حافظه گرافیکی در فرآیند استنباط به شمار می‌رود. پژوهش‌های نوین آشکار ساخته‌اند که تمامی توکن‌ها اهمیت یکسانی در محاسبات توجه ایفا نمی‌کنند؛ بلکه بخش اندکی از توکن‌های آغازین به عنوان چاه‌های توجه (\lr{Attention Sinks}) و برخی توکن‌های کلیدی در طول متن به عنوان ضربه‌زنندگان سنگین (\lr{Heavy Hitters}) بار اصلی توجه را به دوش می‌کشند.

موتور فشرده‌سازی \lr{ZeroKV-Neo Compressor} با تلفیق مدیریت پویای حافظه نهان با پیچیدگی حافظه‌ای محدود $O(1)$ و هرس ساختاریافته ماتریس وزن پروجکشن‌های توجه با اتکا به نُرم $L_2$ طراحی شده است. این ماژول کانال‌های وزن و توکن‌های حافظه نهان را بر اساس ترکیبی از چاه‌های توجه پایدار، پنجره محلی لغزان و پرانرژی‌ترین ابعاد تفکیک کرده و بدون تحمیل افت کیفی فاحش، بار حافظه و محاسبات لایه‌های توجه را کاهش می‌دهد.

\subsection{مبانی نظری و فرمول‌بندی ریاضی}
الگوریتم \lr{ZeroKV-Neo} شامل دو سطح متمایز فشرده‌سازی است: فشرده‌سازی ایستای وزن‌ها و تخلیه پویای حافظه نهان کلید-مقدار.

در سطح فشرده‌سازی وزن لایه‌های پروجکشن، فرض کنید $\mathbf{W} \in \mathbb{R}^{d_{\text{out}} \times d_{\text{in}}}$ ماتریس وزن یک نگاشت توجه باشد. اهمیت هر سطر با نُرم اقلیدسی $L_2$ آن سنجیده می‌شود:
\begin{equation}
n_i = \|\mathbf{w}_i\|_2 = \sqrt{\sum_{j=1}^{d_{\text{in}}} W_{i, j}^2}, \quad \forall i \in \{1, \dots, d_{\text{out}}\}
\end{equation}
با تعیین نسبت نگه‌داری $\alpha \in (0, 1]$، تعداد سطرهای نگه‌داشته‌شده $K = \lfloor d_{\text{out}} \cdot \alpha \rfloor$ تعیین می‌شود. این بودجه به سه گروه تخصیص می‌یابد:
\begin{enumerate}
    \item کانال‌های پیش‌فرض چاه توجه با تعداد $N_{\text{sinks}}$ از سطرهای نخستین: $\mathcal{I}_{\text{sinks}} = \{1, \dots, N_{\text{sinks}}\}$.
    \item کانال‌های ضربه‌زننده سنگین بر اساس بالاترین مقادیر نُرم در میان سطرهای باقی‌مانده:
    \begin{equation}
    \mathcal{I}_{\text{heavy}} = \operatorname{argTop-k}\left(\{n_i \mid i \notin \mathcal{I}_{\text{sinks}}\}, K - N_{\text{sinks}}\right)
    \end{equation}
\end{enumerate}
ماتریس وزن هرس‌شده با ماسک باینری $\mathbf{M}$ حاصل می‌شود:
\begin{equation}
M_i = \begin{cases} 1 & \text{if } i \in \mathcal{I}_{\text{sinks}} \cup \mathcal{I}_{\text{heavy}} \\ 0 & \text{otherwise} \end{cases}, \quad \hat{\mathbf{W}} = \mathbf{W} \odot (\mathbf{M} \mathbf{1}^\top)
\end{equation}

در سطح پویای استنباط (\lr{ZeroKVNeoCache}) با طول حافظه مجاز $L_{\text{max}}$، هنگام فراتر رفتن طول توالی $T > L_{\text{max}}$، توکن‌های نگهداری‌شده از الحاق چاه‌های نخستین، پنجره اخیر و توکن‌های میانی با بیشترین نُرم انتخاب می‌شوند:
\begin{equation}
\mathbf{K}_{\text{bounded}} = \left[\mathbf{K}_{:, :, :N_s, :}; \; \mathbf{K}_{:, :, \mathcal{I}_{\text{mid\_heavy}}, :}; \; \mathbf{K}_{:, :, -N_w:, :}\right]
\end{equation}

\subsection{معماری و پیاده‌سازی}
پیاده‌سازی ماژول در فایل \lr{Codes/framework/compression\_engine/zerokv\_neo\_compressor.py} مستقر بوده و کلاس‌های \lr{ZeroKVNeoCache} و \lr{ZeroKVNeoCompressor} با نام‌های \lr{zerokv}، \lr{zerokv\_neo} و \lr{zerokv\_neo\_compressor} در رجیستری فشرده‌سازی ثبت شده‌اند. شاخص‌ترین مزایای فنی پیاده‌سازی عبارتند از:
\begin{itemize}
    \item \textbf{مدیریت حافظه نهان کراندار (\lr{Bounded $O(1)$ KV Cache}):} کلاس \lr{ZeroKVNeoCache} با مدیریت دقیق ابعاد چهاربعدی تانسورهای توجه $(B, H, N, D)$ مانع از سرریز حافظه در توالی‌های بسیار طولانی می‌گردد.
    \item \textbf{تنظیم نرخ نگه‌داری با تلاش مجدد خودکار:} در صورتی که خطای نسبی فروبنیوس هرس وزن از آستانه مجاز (\lr{max\_error=0.85}) بیشتر شود، الگوریتم در یک حلقه تکرار شونده تا ۳ گام نرخ نگه‌داری را افزایش داده تا خطای بازسازی مهار شود.
    \item \textbf{پشتیبانی اختصاصی از لایه‌های خطی و کانولوشن:} انطباق با نگاشت‌های \lr{Conv1D} در \lr{GPT-2} از طریق ترانهاده‌سازی خودکار پیش از محاسبه نُرم سطرها.
    \item \textbf{تمرکز انحصاری بر پروجکشن‌های توجه (\lr{Target Attention Only}):} اعمال گزینشی هرس تنها بر لایه‌های توجه (\lr{q\_proj, k\_proj, v\_proj, o\_proj}) بدون دستکاری لایه‌های پیش‌خوران به منظور پیشگیری از اعوجاج در بازنمایی واژگانی.
\end{itemize}

\subsection{نتایج تجربی}
نتایج حاصل از ارزیابی‌های جامع بنچ‌مارک فشرده‌ساز \lr{zerokv\_neo\_compressor} روی مدل‌های پایه در جدول \ref{tab:benchmark_zerokv_neo_compressor} درج شده است:

\begin{table}[H]
\centering
\small
\caption{نتایج تجربی و اعتبارسنجی بنچ‌مارک فشرده‌ساز \lr{zerokv\_neo\_compressor} روی مدل‌های پایه}
\label{tab:benchmark_zerokv_neo_compressor}
\begin{tabular}{lccccc}
\toprule
\textbf{مدل پایه} \lr{(Base Model)} & \textbf{دقت} \lr{(Accuracy)} & \textbf{شاخص} \lr{(F1-Score)} & \textbf{حافظه} \lr{(Peak VRAM)} & \textbf{توان} \lr{(Throughput)} & \textbf{وضعیت} \\
\midrule
\lr{DistilBERT} & \lr{75.92\%} & \lr{0.7925} & \lr{312.19 MB} & \lr{52,122.59 tok/s} & \lr{PASS} \\
\lr{RoBERTa} & \lr{50.92\%} & \lr{0.6748} & \lr{531.69 MB} & \lr{23,289.24 tok/s} & \lr{PASS\_LOW} \\
\lr{GPT-2} & \lr{51.15\%} & \lr{0.6758} & \lr{413.50 MB} & \lr{14,040.51 tok/s} & \lr{PASS\_LOW} \\
\lr{Qwen} & \lr{---} & \lr{---} & \lr{---} & \lr{---} & \lr{NOT\_RUN} \\
\bottomrule
\end{tabular}
\end{table}

بر پایه نتایج تجربی، در مدل \lr{DistilBERT} فشرده‌ساز \lr{zerokv\_neo\_compressor} با ثبت دقت \lr{75.92\%}، شاخص F1 معادل \lr{0.7925} و مصرف حافظه بسیار پایین \lr{312.19 MB} وضعیت موفقیت‌آمیز \lr{PASS} را احراز نموده است. در مدل‌های خودبازگشت و رمزگذار عمیق نظیر \lr{RoBERTa} و \lr{GPT-2}، به دلیل صفر کردن مستقیم سطرهای وزن بدون بازآموزی جبرانی، دقت مدل به حدود \lr{51\%} تنزل یافته و وضعیت \lr{PASS\_LOW} به ثبت رسیده است. این ماژول در مدل \lr{Qwen} آزمایش نشده و وضعیت \lr{NOT\_RUN} دارد.

\subsection{محدودیت‌ها و تحلیل شکست}
تحلیل عوامل افت کارایی و محدودیت‌های الگوریتم حاکی از نکات زیر است:
\begin{itemize}
    \item \textbf{افت ناشی از صفر کردن سخت سطرها بدون بازآموزی (\lr{Hard Zero Pruning}):} صفر کردن وزن‌های سطرها بدون آموزش جبرانی موجب حذف برخی سیگنال‌های پروجکشن سرها شده و در مدل‌های بزرگ‌تری که حساسیت بالا به هماهنگی سرها دارند، دقت اولیه را کاهش می‌دهد.
    \item \textbf{انحراف در فشرده‌سازی لایه‌های دوطرفه:} در مدل‌های دوطرفه (\lr{Bidirectional Encoder}) مانند \lr{RoBERTa}، مفهوم چاه‌های توجه با مدل‌های زبانی خودبازگشت متفاوت است؛ چرا که توکن‌های پایانی نیز به صورت یکنواخت به توکن‌های پیشین متصل می‌شوند و اتکای صِرف به چاه‌های اولیه موجب انحراف میدان توجه می‌گردد.
    \item \textbf{وابستگی دقت به آستانه خطای فروبنیوس:} در صورتی که آستانه خطای مجاز بیش از اندازه آزاد در نظر گرفته شود (مانند $0.85$)، هرس شدید رخ می‌دهد و بازیابی دقت نیازمند استفاده از مکانیزم بازگشت به وزن‌های دست‌نخورده لایه خواهد بود.
\end{itemize}
